\thispagestyle{empty}

% Abstract y resumen
\abstract{Resumen}{Presentamos en estas notas de clase algunos problemas rutinarios sobre integración para el curso Cálculo Multivariado. En ellas se trabaja, básicamente, la integración de una función sobre un objeto que puede ser una curva, superficie o sólido, con el fin de parametrizar la región de \mbox{integración}. Esta idea es diferenciadora con respecto a la mayoría de textos de cálculo multivariado. Cada problema ofrece un gráfico asociado en diversas vistas y el cálculo que se requiere. Además se incluyen las instrucciones utilizadas en el programa \verb"Wolfram Mathematica". El cambio de sistema de coordenadas para evaluar integrales triples se hace tradicionalmente utilizando coordenadas esféricas y cilíndricas.}{Palabras clave}{integración múltiple, parametrización, campos escalares, campos vectoriales, Wolfram Mathematica.}

\rule{3cm}{0.5pt}\vspace{\parskip}

\begin{otherlanguage}{english}
\abstract{Abstract}{In these lecture notes, we present some routine integration problems for the Multivariable Calculus course. They primarily focus on integrating a function over an object—which may be a curve, a surface, or a solid—in order to parameterize the region of integration.
This approach sets it apart from most multivariable calculus textbooks. Each problem includes an associated graph from various views and the required calculation. Additionally, the commands used in the \verb"Wolfram Mathematica"  program are included. The change of coordinate system for evaluating triple integrals is traditionally done using spherical and cylindrical coordinates.}{Keywords}{multiple integration, parameterization, scalar fields, vector fields, Wolfram \mbox{Mathematica}.}
\end{otherlanguage}

\vfill

\hfill\comocitar{Zambrano Caviedes, J. y Baquero Torres, E. (2026). \emph{Problemas resueltos de integración múltiple con aplicaciones en Wolfram Mathematica.} Universidad Distrital Francisco José de Caldas-Editorial~UD.}{10.69740/UD.9786287926035.9786287926059}

